\section{Preliminaries and Problem Formulation}
\label{sec:problem}

We study how many input tuples must be deleted to move an aggregate-query
answer across a prescribed threshold. This section specifies the query
model and defines budgeted instability witnesses. Section~\ref{sec:complexity}
studies their complexity; the subsequent coverage analysis is reserved
for Section~\ref{sec:theory}.

\subsection{Data and Query Model}
\label{sec:setting}

Let $D$ be a finite relational database instance. We identify each tuple
occurrence by a distinct identifier, so equal-valued tuples can be deleted
independently. A query $Q$ consists of a positive select--project--join
(SPJ) expression followed by a scalar aggregate. We use bag semantics:
different successful derivations contribute separately, including when
projection produces equal-valued results. We consider \agg{COUNT},
\agg{SUM}, \agg{MAX}, \agg{MIN}, and \agg{AVG}, without grouping,
negation, or duplicate elimination. For \agg{SUM}, contributions are
nonnegative. Aggregate inputs are numeric and contain no null values.

Only tuples in a designated set $\Del\subseteq D$ may be deleted; all other
tuples are immutable. A deletion set $S\subseteq\Del$ produces the instance
$D\setminus S$. Deletions can remove derivations but cannot create new ones.
There is no restriction on overlap between derivations or on which input
relations contain deletable tuples.

Let $\B$ be the multiset of successful derivations before aggregation.
For a derivation $b\in\B$, let $v_b$ be its aggregate input value and let
$T_b\subseteq\Del$ contain the deletable tuples on which it depends.
Immutable tuples are implicit in the derivation and do not appear in $T_b$.
The surviving derivations after deleting $S$ are
\begin{equation}
\label{eq:surviving-derivations}
\B(S)=\{b\in\B:S\cap T_b=\emptyset\},
\end{equation}
where multiplicities are retained. Thus, $b$ is destroyed exactly when
$S\cap T_b\ne\emptyset$. This representation follows the provenance view of query
evaluation~\cite{buneman2001whywhere,green2007provenance,amsterdamer2011aggregate}.
It describes the query semantics; it does not require an algorithm to
materialize every derivation.

We define \agg{COUNT} and \agg{SUM} on an empty input to be zero,
using the additive identity for \agg{SUM} rather than SQL's null result.
For \agg{MAX}, \agg{MIN}, and \agg{AVG}, a valid deletion set must
leave $\B(S)\ne\emptyset$. We assume that the original aggregate is defined. If a tuple is used in
multiple atom positions of a self-join derivation, it occurs only once in
$T_b$: deleting that tuple destroys the derivation once.

\paragraph{Computational input.}
For algorithmic claims, values and thresholds are binary-encoded rationals,
and the deletion budget $k$ is a binary-encoded nonnegative integer.
Data-complexity statements fix $Q$; combined complexity allows it to vary.
Write $n=|\Del|$, $M=|\B|$, and $L=\sum_{b\in\B}|T_b|$.
An explicit-derivation input lists every value and support, including bag
multiplicities. Alternatively, a database engine can evaluate residual
queries, or an explicitly specified circuit can provide value oracles.
These representations have different costs; their execution methods are
reserved for Section~\ref{sec:algorithms}.
For a fixed query with $a$ relation-atom occurrences, $M\le |D|^a$,
so enumeration is polynomial in data size, though it may be impractical.
This is not a polynomial bound in combined query/database size or in the
size of an arbitrary compressed representation. Arithmetic-operation bounds
below are distinguished from bit complexity; exact rational operations
also depend on operand lengths.

\subsection{Instability Witnesses}
\label{sec:witnesses}

Fix a deletion budget $k\in\mathbb{Z}_{\ge0}$ and a required answer
change $\theta>0$. The budget counts source tuple occurrences; the
threshold measures the change in the aggregate answer. For \agg{COUNT}, \agg{SUM},
\agg{MAX}, and \agg{AVG}, we seek a downward change and write
$\tau_{\downarrow}=Q(D)-\theta$. For \agg{MIN}, deletions cannot decrease
the minimum of a nonempty input, so we instead seek an upward change and
write $\tau_{\uparrow}=Q(D)+\theta$. Accordingly, define
\begin{equation}
\label{eq:deletion-benefit}
\DeltaQ(S)=
\begin{cases}
Q(D\setminus S)-Q(D), & \text{for \agg{MIN}},\\
Q(D)-Q(D\setminus S), & \text{otherwise}.
\end{cases}
\end{equation}
The threshold is non-strict: equality suffices to establish instability.
Unlike the other operators, \agg{AVG} can move in either direction under
deletions. We state its downward version; negating all values converts an
upward instance into a downward one with the same witnesses. A two-sided
\agg{AVG} request is solved by taking the smaller feasible witness from
the two directions, reporting infeasibility only if both fail.

\begin{definition}[$k$-instability witness]
\label{def:witness}
For a prescribed budget $k$ and threshold $\theta$, a valid deletion set
$S\subseteq\Del$ is a \emph{$k$-instability witness} if
\[
|S|\le k\qquad\text{and}\qquad\DeltaQ(S)\ge\theta.
\]
Equivalently, within the deletion budget it satisfies
$Q(D\setminus S)\le\tau_{\downarrow}$ for a downward threshold, or
$Q(D\setminus S)\ge\tau_{\uparrow}$ for \agg{MIN}.
\end{definition}

\begin{problem}[Minimum $k$-Instability Witness (\MIW)]
\label{prob:miw}
Given $D$, $Q$, $\Del$, $k\ge0$, and $\theta>0$, compute
\begin{equation}
\label{eq:miw}
\begin{aligned}
S^\star\in\argmin_{S\subseteq\Del}\{\,|S|:\;&S\text{ is valid},\\
&|S|\le k,\ \DeltaQ(S)\ge\theta\,\},
\end{aligned}
\end{equation}
or report that no witness exists within budget $k$.
\end{problem}

We require minimum cardinality, rather than merely inclusion minimality.
All deletions have unit cost. The parameters have different roles: $k$
bounds the allowed intervention, while $\theta$ fixes the required answer
change. Minimization chooses the smallest witness among those satisfying
both constraints; the budget is an upper bound, not a required size.

Let $\rho_\theta(D,Q,\Del)$ denote the minimum valid witness size without
a budget, taking value $\infty$ when no deletion set reaches the threshold.
Then a $k$-instability witness exists exactly when
\[
\rho_\theta(D,Q,\Del)\le k.
\]
When feasible, Problem~\ref{prob:miw} returns this minimum size. Failure
within a particular budget means $\rho_\theta>k$; it does not by itself
establish that the threshold is unreachable with more deletions. A larger
$\rho_\theta$ means greater directional threshold stability. In particular,
$k=0$ admits no witness because $\theta>0$.

\begin{example}
\label{ex:miw}
Consider the join $R(K)\Join V(K,A)$, with deletable tuples
$r_1=R(1)$ and $r_2=R(2)$, and immutable tuples
$V(1,4)$, $V(1,6)$, and $V(2,5)$. The sum of $A$ is $15$.
Deleting $r_1$ destroys the two derivations with values $4$ and $6$,
whereas deleting $r_2$ destroys the derivation with value $5$.
For $\theta=8$ and $k=1$, the set $\{r_1\}$ is a minimum
$1$-instability witness:
the remaining sum is $5\le15-8$, while $\{r_2\}$ does not reach the
threshold. The deletion budget concerns input tuples, not the number of
destroyed join results.
\end{example}

\paragraph{Aggregate-specific structure.}
For \agg{COUNT} and nonnegative \agg{SUM}, the numerical change in
the answer is a weighted-coverage function over destroyed derivations.
For \agg{MAX} and \agg{MIN}, the threshold instead identifies a set
of critical derivations that must all be destroyed, while at least one
other derivation survives. For \agg{AVG}, the numerical change need
not be monotone or submodular. These distinctions organize the planned
coverage and approximation analysis in Section~\ref{sec:theory}.

\paragraph{Why bag projection is essential.}
Let deletable $R(1),R(2)$ join with immutable $V(1,7),V(2,7)$ and project
$A$. Under bag semantics the result is $\{7,7\}$, and deleting either
$R$ tuple removes one copy. Duplicate elimination instead gives $\{7\}$: either deletion alone
leaves that value present, while both remove it.
The marginal \agg{COUNT} loss of the second deletion rises from zero to
one. Thus duplicate elimination can invalidate submodularity; results for
set projection cannot be imported without checking their semantics.
